%! Adding \& Subtracting Rational Expressions — Unit 4, Lesson 4.3 — Warm-Up
\documentclass[10pt]{article}
\usepackage{algebra2-article}
\usepackage{algebra2-boxes}
\newcommand{\TallMath}[1]{$\displaystyle #1\rule[-1.4em]{0pt}{3.2em}$}

\begin{document}

\pageheader{Unit 4, Lesson 4.3}{Warm-Up}
\namedateperiod

\vspace{0.10in}

\begin{notesbox}{Getting Ready: Skills We'll Use Today}
\small Three skills, one per leg of today's lesson. Work quickly.

\vspace{4pt}
\textbf{1. Number fractions --- the common denominator.}
\begin{enumerate}[label=(\alph*), leftmargin=2.1em, itemsep=3pt]
  \item Write each denominator as a product of primes, then build the \emph{smallest} denominator
        that both can become. \\[2pt]
        $12=\blank{1.4cm}$, \quad $18=\blank{1.4cm}$, \quad so the least common denominator is
        \blank{1.2cm}. \\[3pt]
        $\dfrac{5}{12}+\dfrac{7}{18}=\blank{1.2cm}+\blank{1.2cm}=\blank{1.4cm}$
  \item You could also have used $12\cdot18=216$ as the denominator. Rewrite both fractions over
        $216$ and add: \\[3pt]
        \blank{1.6cm} $+$ \blank{1.6cm} $=$ \blank{1.8cm} \\[3pt]
        Same value as part (a)? \blank{1.0cm} \quad What extra work did it cost you?
        \blank{4.2cm}
  \item To \emph{divide} by a fraction, multiply by its \blank{2.2cm} (Lesson 4.2):
        $\dfrac{2}{3}\div\dfrac{5}{6}=\dfrac{2}{3}\cdot\blank{1.0cm}=\blank{1.0cm}$
\end{enumerate}

\vspace{4pt}
\textbf{2. Signs and factors.}
\begin{enumerate}[label=(\alph*), leftmargin=2.1em, itemsep=3pt]
  \item $9-(4-6)=\blank{1.0cm}$ \qquad but \qquad $9-4-6=\blank{1.0cm}$. \\[2pt]
        They are not equal, because a minus sign in front of a parenthesis changes
        \blank{3.6cm} inside it.
  \item Factor completely (the Unit 3 toolkit): \; $x^2-9=\blank{3.0cm}$; \quad
        $x^2-4=\blank{3.0cm}$; \quad $x^2+5x+6=\blank{3.0cm}$
\end{enumerate}

\vspace{4pt}
\textbf{3. Simplify and restrict} (Lesson 4.1). For $\dfrac{3x-9}{x^2-9}$:
\begin{enumerate}[label=(\alph*), leftmargin=2.1em, itemsep=3pt]
  \item Factored form: \blank{3.6cm} \quad Restrictions: $x\neq$ \blank{1.8cm}
  \item Simplest form: \blank{2.0cm}
  \item Which excluded value is a \emph{hole}? \blank{1.2cm} \; Which is a \emph{wall} (vertical
        asymptote)? \blank{1.2cm}
\end{enumerate}

\end{notesbox}

\end{document}
